\silentchapter{Omhugsmal}{omhugsmal}
\noindent
\begin{widebody}
\scriptsize
\setlength{\tabcolsep}{2pt}
\renewcommand{\arraystretch}{1.25}
\noindent\begin{tabular}{@{}rp{20mm}p{24mm}p{16mm}p{30mm}@{}}
\toprule
\textbf{\#} & \textbf{Lag} & \textbf{Problemrom} & \textbf{Tidsrom} & \textbf{Omhug-krav} \\
\midrule
17 & Noosfære, sivilisasjon & Idérom, kunnskapssystem, kollektive prosjekt & Hundreår til millennium & Vitskap, kunst, lovverk, planetær koordinering; reproduserer verdsmodellar \\
16 & Samfunn, institusjon & Sosialt rom, økonomi, politikk, jus & Tiår til hundreår & Arbeidsdeling, rettsstat, utdanning; reproduserer seg over generasjonar \\
15 & Gruppe, familie & Relasjonsrom, nære band, delt prosjekt & År til tiår & Omsorg, oppseding, arbeidsfellesskap; byggjer tillit over tid \\
14 & Person, sjølv & Biografisk rom, livsprosjekt, identitet & Tiår + forestilt rom & Autobiografisk minne, viljestyrt handling, symbolsk tenking, moralsk agens \\
13 & Hjerne, sentralnervesystem & Åtferdsrom, persepsjon, planar & Millisekund til år & Persepsjon, minne, prediksjon, læring, språk, emosjon \\
12 & Immunsystem & Sjølv, ikkje-sjølv, molekylær identitet & Sekund til livstid & Skil sjølv frå ikkje-sjølv; tolererer eller angrip; regulerer nervesystem \\
11 & Organsystem & Funksjonelt rom, homeostase, allostase & Sekund til døgn & Koordinert funksjon: pumping, filtrering, fordøying, pust, forsvar \\
10 & Organ & Fysiologisk rom, spesialisert organfunksjon & Sekund til timar & Utfører éi hovudoppgåve: hjarte pumpar, lever detoksifiserer, nyre filtrerer \\
9 & Vev & Morfogenetisk rom, anatomisk form & Minutt til dagar & Morfogenese, regenerering, formminne; «kor er eg i kroppen» \\
8 & Celle & Beslutningsrom: del, dø, migrer, differensier & Sekund til timar & Metabolisme, sansing, rørsle, kommunikasjon, basal kognisjon \\
7 & Organell & Funksjonelt subrom i cella (mitokondrium, kjerne) & Millisek.\ til minutt & Energiproduksjon, proteinsyntese, informasjonslagring \\
6 & Gen-regulatorisk nettverk & Uttrykksrom: kva gen er av, på & Minutt til timar & Sensing av kontekst; seks pavlovske læringsformer \\
5 & Enzym, proteinkompleks & Konformasjonsrom, aktivt sete, allosteriske setar & Mikrosek.\ til sekund & Katalyse: endrar kjemisk identitet til substrat; primitivt minne (hysterese) \\
4 & Molekyl & Bindingsrom, van der Waals-overflate, elektronsky & Nano- til mikrosekund & Stabil struktur, reaktivitet, polaritet; inngår i større kompleks \\
3 & Atom & Elektronrom: skal, bindingar, ladning & Femto- til nanosekund & Danne bindingar, ionisere, emittere eller absorbere foton \\
2 & Subatomær partikkel & Spinn, ladning, kvantetal & Atto- til femtosekund & Bere ladning, spinn, magnetisk moment; inngå i atom \\
1 & Kvark, leptonfelt & Kvantetilstandsrom, farge, smak & Under attosekund & Interagere via fargekraft; konstituere nukleonar \\
\bottomrule
\end{tabular}
\end{widebody}
